\begin{examplebox}
	\textbf{\textcolor{CExample}{例1｜基础：幂除以对数}}
	\textbf{题目：}$\displaystyle\lim_{x\to+\infty}\frac{x^3}{\log_2x}$。

	\textbf{最快思路：}分子是$x^3$，分母是$\log_2x$；底数$2>1$，幂指数$3>0$。\textbf{幂最终比对数增长快}，所以
	\[
		\boxed{\lim_{x\to+\infty}\frac{x^3}{\log_2x}=+\infty}.
	\]
	\textbf{用数字感受：}当$x=2,4,8$时，比值分别是$8,32,\frac{512}{3}$，不断变大。数字只帮助理解，结论来自上面的证明。
	\begin{center}
	\begin{tikzpicture}[>=stealth,font=\small]
		\node[draw=blue!60,rounded corners=3pt,inner sep=5pt] (a) at (0,0) {$x^3$：增长更快};
		\node[draw=green!55!black,rounded corners=3pt,inner sep=5pt] (b) at (5.8,0) {$\log_2x$：增长较慢};
		\draw[->,very thick,blue!75!black] (a.east)-- (b.west);
	\node[blue!75!black] at (2.9,1) {相差倍数越来越大};
\end{tikzpicture}

	\end{center}
	\textbf{记住：}较快的在分子，结果趋向$+\infty$。

	\medskip
	\textbf{\textcolor{CExample}{例2｜变形：幂除以指数}}
	\textbf{题目：}$\displaystyle\lim_{x\to+\infty}\frac{x^{100}}{2^x}$。

	\textbf{最快思路：}虽然$100$很大，但它是\textbf{固定的指数}；$2^x$才是增长更快的那一个。
	\[
		\frac{2^x}{x^{100}}\to+\infty
		\quad\Longrightarrow\quad
		\boxed{\lim_{x\to+\infty}\frac{x^{100}}{2^x}=0}.
	\]
	\begin{center}
	\begin{tikzpicture}[>=stealth,font=\small]
		\node[draw=red!70!black,rounded corners=3pt,inner sep=5pt] (a) at (0,0) {$\dfrac{x^{100}}{2^x}$};
		\node[draw=gray!70,rounded corners=3pt,inner sep=5pt] (b) at (4.7,0) {分母增长更快 $\Rightarrow0$};
		\draw[->,thick] (a.east)--(b.west);
	\end{tikzpicture}

	\end{center}
	\textbf{易忽略：}$x$不太大时，$x^{100}$可能更大；题目问的是$x\to+\infty$。

	\medskip
	\textbf{\textcolor{CExample}{例3｜综合：分子有两种函数}}
	\textbf{题目：}$\displaystyle\lim_{x\to+\infty}\frac{x^{10}+\log_3x}{4^x}$。

	\textbf{第一步：}分子中$\log_3x$比$x^{10}$增长慢，先把$x^{10}$提出来：
	\[
		\frac{x^{10}+\log_3x}{4^x}
		=\underbrace{\frac{x^{10}}{4^x}}_{\to0}
		\left(1+\underbrace{\frac{\log_3x}{x^{10}}}_{\to0}\right).
	\]
	\textbf{第二步：}括号趋于$1$，前面的分式趋于$0$，所以
	\[
		\boxed{\lim_{x\to+\infty}\frac{x^{10}+\log_3x}{4^x}=0}.
	\]
	\begin{center}
	\begin{tikzpicture}[>=stealth,font=\small]
		\node[draw=red!70!black,rounded corners=3pt,inner sep=4pt] (p) at (0,0) {$x^{10}$};
		\node[draw=green!50!black,rounded corners=3pt,inner sep=4pt] (l) at (2.6,0) {$\log_3x$};
		\node[draw=blue!65!black,rounded corners=3pt,inner sep=4pt] (e) at (5.7,0) {$4^x$};
		\draw[->,thick] (l.south)--++(0,-0.65) node[below] {比分子第一项慢};
		\draw[->,thick] (p.north) to[bend left=25] node[above] {最终被超过} (e.north);
	\end{tikzpicture}

	\end{center}
	\textbf{记住：}先找分子中增长最快的项，再和分母比较。这里用乘法写出全过程，避免“凭感觉直接删项”。
\end{examplebox}
